% Preámbulo para compilar doble_ponderacion.md a PDF con XeTeX (Tectonic).
%
% Resuelve los 312 caracteres que la fuente por defecto (Latin Modern) no tiene
% y que salían como hueco en blanco:
%
%   ─ │ ├ └ ┐ ┘   (144+60+38+18+2+2) el diagrama de la sección 2
%   ✓             (34) la tabla de recuento de la sección 2.1
%   α ρ           (10+4) el parámetro de la sección 5.1 y la rho de Spearman
%
% Dos estrategias distintas a propósito:
%
%  - Los símbolos matemáticos (α ρ ✓) se mapean a comandos de LaTeX, que
%    componen igual en cualquier equipo y no dependen de ninguna fuente
%    instalada.
%  - Los de dibujo de caja NO tienen equivalente en LaTeX, así que hacen falta
%    de una fuente real. Consolas viene con Windows desde Vista y los trae
%    todos; se verificó que existe antes de fijarla aquí.

\usepackage{newunicodechar}
\newunicodechar{α}{\ensuremath{\alpha}}
\newunicodechar{ρ}{\ensuremath{\rho}}
\newunicodechar{✓}{\ensuremath{\checkmark}}

% Consolas sólo para el texto monoespaciado: es donde vive el diagrama. El
% cuerpo del documento sigue con la serif de LaTeX, que compone mucho mejor.
\setmonofont{Consolas}[Scale=0.82]

% El diagrama de la sección 2 mide 74 columnas y se sale de la caja con el
% cuerpo por defecto. Reducirlo aquí evita tener que estrechar el margen de
% todo el documento por culpa de un solo bloque.
\usepackage{fancyvrb}
\fvset{fontsize=\footnotesize}

% Las 18 tablas son anchas; sin esto el texto de las celdas queda con espacios
% enormes entre palabras.
\usepackage{ragged2e}
\usepackage{array}

% Deja partir las palabras en monoespaciada. Sin esto, el nombre del PDF de
% origen --63 caracteres sin espacios-- no cabe en la caja y se sale 18,45 pt
% por el margen derecho, que es visible en la página 1.
\usepackage[htt]{hyphenat}

% Filetes de tabla con aire, que es lo que separa una tabla legible de una
% rejilla apretada.
\usepackage{booktabs}
\setlength{\aboverulesep}{2pt}
\setlength{\belowrulesep}{2pt}
\renewcommand{\arraystretch}{1.15}
